
\documentclass{article}
\usepackage{colortbl}
\usepackage{makecell}
\usepackage{multirow}
\usepackage{supertabular}

\begin{document}

\newcounter{utterance}

\twocolumn

{ \footnotesize  \setcounter{utterance}{1}
\setlength{\tabcolsep}{0pt}
\begin{supertabular}{c@{$\;$}|p{.15\linewidth}@{}p{.15\linewidth}p{.15\linewidth}p{.15\linewidth}p{.15\linewidth}p{.15\linewidth}}

    \# & $\;$A & \multicolumn{4}{c}{Game Master} & $\;\:$B\\
    \hline 

    \theutterance \stepcounter{utterance}  

    & & \multicolumn{4}{p{0.6\linewidth}}{\cellcolor[rgb]{0.9,0.9,0.9}{%
	\makecell[{{p{\linewidth}}}]{% 
	  \tt {\tiny [A$\langle$GM]}  
	 Du nimmst an einem kollaborativen Verhandlungspiel Teil.\\ \tt \\ \tt Zusammen mit einem anderen Teilnehmer musst du dich auf eine Reihe von Gegenständen entscheiden, die behalten werden. Jeder von euch hat eine persönliche Verteilung über die Wichtigkeit der einzelnen Gegenstände. Jeder von euch hat eine eigene Meinung darüber, wie wichtig jeder einzelne Gegenstand ist (Gegenstandswichtigkeit). Du kennst die Wichtigkeitsverteilung des anderen Spielers nicht. Zusätzlich siehst du, wie viel Aufwand jeder Gegenstand verursacht.  \\ \tt Ihr dürft euch nur auf eine Reihe von Gegenständen einigen, wenn der Gesamtaufwand der ausgewählten Gegenstände den Maximalaufwand nicht überschreitet:\\ \tt \\ \tt Maximalaufwand = 4026\\ \tt \\ \tt Hier sind die einzelnen Aufwände der Gegenstände:\\ \tt \\ \tt Aufwand der Gegenstände = {"C33": 780, "A89": 742, "A48": 997, "A04": 665, "A41": 140, "C64": 77, "C49": 513, "A61": 383, "C04": 587, "B53": 825, "B85": 319, "B61": 448, "A45": 516, "A03": 694, "B21": 366}\\ \tt \\ \tt Hier ist deine persönliche Verteilung der Wichtigkeit der einzelnen Gegenstände:\\ \tt \\ \tt Werte der Gegenstandswichtigkeit = {"C33": 138, "A89": 583, "A48": 868, "A04": 822, "A41": 783, "C64": 65, "C49": 262, "A61": 121, "C04": 508, "B53": 780, "B85": 461, "B61": 484, "A45": 668, "A03": 389, "B21": 808}\\ \tt \\ \tt Ziel:\\ \tt \\ \tt Dein Ziel ist es, eine Reihe von Gegenständen auszuhandeln, die dir möglichst viel bringt (d. h. Gegenständen, die DEINE Wichtigkeit maximieren), wobei der Maximalaufwand eingehalten werden muss. Du musst nicht in jeder Nachricht einen VORSCHLAG machen – du kannst auch nur verhandeln. Alle Taktiken sind erlaubt!\\ \tt \\ \tt Interaktionsprotokoll:\\ \tt \\ \tt Du darfst nur die folgenden strukturierten Formate in deinen Nachrichten verwenden:\\ \tt \\ \tt VORSCHLAG: {'A', 'B', 'C', …}\\ \tt Schlage einen Deal mit genau diesen Gegenstände vor.\\ \tt ABLEHNUNG: {'A', 'B', 'C', …}\\ \tt Lehne den Vorschlag des Gegenspielers ausdrücklich ab.\\ \tt ARGUMENT: {'...'}\\ \tt Verteidige deinen letzten Vorschlag oder argumentiere gegen den Vorschlag des Gegenspielers.\\ \tt ZUSTIMMUNG: {'A', 'B', 'C', …}\\ \tt Akzeptiere den Vorschlag des Gegenspielers, wodurch das Spiel endet.\\ \tt STRATEGISCHE ÜBERLEGUNGEN: {'...'}\\ \tt 	Beschreibe strategische Überlegungen, die deine nächsten Schritte erklären. Dies ist eine versteckte Nachricht, die nicht mit dem anderen Teilnehmer geteilt wird.\\ \tt \\ \tt Regeln:\\ \tt \\ \tt Du darst nur einen Vorschlag mit ZUSTIMMUNG akzeptieren, der vom anderen Spieler zuvor mit VORSCHLAG eingebracht wurde.\\ \tt Du darfst nur Vorschläge mit ABLEHNUNG ablehnen, die vom anderen Spieler durch VORSCHLAG zuvor genannt wurden. \\ \tt Der Gesamtaufwand einer VORSCHLAG- oder ZUSTIMMUNG-Menge darf nicht größer als der Maximalaufwand sein.  \\ \tt Offenbare deine versteckte Wichtigkeitsverteilung nicht.\\ \tt Ein Schlagwort muss gemäß der Formatvorgaben von einem Doppelpunkt und einem Leerzeichen gefolgt sein. Das Argument muss eine Python-Menge mit 0 oder mehr Strings sein.  \\ \tt Die Form muss also SCHLAGWORT: {...} sein.\\ \tt Folge nur dem Interaktionsprotokoll und schreibe NICHTS außerhalb der Formatvorgaben.\\ \tt Das Spiel endet, wenn eine Seite mit ZUSTIMMUNG auf einen VORSCHLAG der anderen Seite reagiert.  \\ \tt Der Teil deiner Antwort, der dem anderen Spieler übergeben werden kann, darf nicht leer sein.  \\ \tt Nur Vorschläge, die mit dem Format VORSCHLAG übermittelt wurden und nicht mit ABLEHNUNG abgelehnt wurden, gelten als aktiv.  \\ \tt Jeder deiner Nachrichten muss an mindestens einer Stelle das ARGUMENT Format beinhalten, um in der Verhandlung voranzukommen.\\ \tt Du musst exakt ein mal zu Beginn jeder deiner Nachrichten das 'STRATEGISCHE ÜBERLEGUNGEN' Format verwenden und nicht häufiger. Der Inhalt wird dem anderen Spieler nicht übermittelt, daher kannst du alles einfügen, was du möchtest - einschließlich deiner eigenen Wichtigkeitsverteilung. Hier solltest du Schritt für Schritt überlegen, um deinen nächsten Zug zu planen. 
	  } 
	   } 
	   } 
	 & \\ 
 

    \theutterance \stepcounter{utterance}  

    & \multicolumn{4}{p{0.6\linewidth}}{\cellcolor[rgb]{0.8,1,0.9}{%
	\makecell[{{p{\linewidth}}}]{% 
	  \tt {\tiny [A$\rangle$GM]}  
	 assistant\\ \tt STRATEGISCHE ÜBERLEGUNGEN: {'Optimale Aufstellung meinerseits: A48, A04, B21, A41, C04, B61, B85. Gesamtaufwand 3964. Nutzen 4734. Ziel: maximieren meiner Wertigkeit, Rest für Verhandlung.'}\\ \tt VORSCHLAG: {'A48', 'A04', 'B21', 'A41', 'C04', 'B61', 'B85'}\\ \tt ARGUMENT: {'Diese Kombination maximiert meinen persönlichen Nutzen und liegt deutlich unter dem Maximalaufwand. Ich schlage vor, dass wir diese Gegenstände als Ausgangspunkt für eine einvernehmliche Lösung nutzen.'} 
	  } 
	   } 
	   } 
	 & & \\ 
 

    \theutterance \stepcounter{utterance}  

    & & \multicolumn{4}{p{0.6\linewidth}}{\cellcolor[rgb]{0.9,0.9,0.9}{%
	\makecell[{{p{\linewidth}}}]{% 
	  \tt {\tiny [A$\langle$GM]}  
	 Deine Antwort entsprach nicht mit dem korrekten Format für strategische Überlegungen entsprochen. Das erste strukturierte Format muss folgende Form haben: STRATEGISCHE ÜBERLEGUNGEN: {...}. Versuche es erneut. 
	  } 
	   } 
	   } 
	 & \\ 
 

    \theutterance \stepcounter{utterance}  

    & \multicolumn{4}{p{0.6\linewidth}}{\cellcolor[rgb]{0.8,1,0.9}{%
	\makecell[{{p{\linewidth}}}]{% 
	  \tt {\tiny [A$\rangle$GM]}  
	 STRATEGISCHE ÜBERLEGUNGEN: {"Optimale Aufstellung meinerseits: A48, A04, B21, A41, C04, B61, B85. Gesamtaufwand 3964. Nutzen 4734. Ziel: maximieren meiner Wertigkeit. Rest für Verhandlung."}\\ \tt VORSCHLAG: {'A48', 'A04', 'B21', 'A41', 'C04', 'B61', 'B85'}\\ \tt ARGUMENT: {'Diese Kombination maximiert meinen persönlichen Nutzen und liegt deutlich unter dem Maximalaufwand. Ich schlage vor, dass wir diese Gegenstände als Ausgangspunkt für eine einvernehmliche Lösung nutzen.'} 
	  } 
	   } 
	   } 
	 & & \\ 
 

    \theutterance \stepcounter{utterance}  

    & & \multicolumn{4}{p{0.6\linewidth}}{\cellcolor[rgb]{0.9,0.9,0.9}{%
	\makecell[{{p{\linewidth}}}]{% 
	  \tt {\tiny [GM$\rangle$B]}  
	 Du nimmst an einem kollaborativen Verhandlungspiel Teil.\\ \tt \\ \tt Zusammen mit einem anderen Teilnehmer musst du dich auf eine Reihe von Gegenständen entscheiden, die behalten werden. Jeder von euch hat eine persönliche Verteilung über die Wichtigkeit der einzelnen Gegenstände. Jeder von euch hat eine eigene Meinung darüber, wie wichtig jeder einzelne Gegenstand ist (Gegenstandswichtigkeit). Du kennst die Wichtigkeitsverteilung des anderen Spielers nicht. Zusätzlich siehst du, wie viel Aufwand jeder Gegenstand verursacht.  \\ \tt Ihr dürft euch nur auf eine Reihe von Gegenständen einigen, wenn der Gesamtaufwand der ausgewählten Gegenstände den Maximalaufwand nicht überschreitet:\\ \tt \\ \tt Maximalaufwand = 4026\\ \tt \\ \tt Hier sind die einzelnen Aufwände der Gegenstände:\\ \tt \\ \tt Aufwand der Gegenstände = {"C33": 780, "A89": 742, "A48": 997, "A04": 665, "A41": 140, "C64": 77, "C49": 513, "A61": 383, "C04": 587, "B53": 825, "B85": 319, "B61": 448, "A45": 516, "A03": 694, "B21": 366}\\ \tt \\ \tt Hier ist deine persönliche Verteilung der Wichtigkeit der einzelnen Gegenstände:\\ \tt \\ \tt Werte der Gegenstandswichtigkeit = {"A48": 65, "A04": 121, "B21": 138, "A41": 262, "B53": 389, "A45": 461, "A89": 484, "C04": 508, "B61": 583, "B85": 668, "A03": 780, "C49": 783, "C33": 808, "A61": 822, "C64": 868}\\ \tt \\ \tt Ziel:\\ \tt \\ \tt Dein Ziel ist es, eine Reihe von Gegenständen auszuhandeln, die dir möglichst viel bringt (d. h. Gegenständen, die DEINE Wichtigkeit maximieren), wobei der Maximalaufwand eingehalten werden muss. Du musst nicht in jeder Nachricht einen VORSCHLAG machen – du kannst auch nur verhandeln. Alle Taktiken sind erlaubt!\\ \tt \\ \tt Interaktionsprotokoll:\\ \tt \\ \tt Du darfst nur die folgenden strukturierten Formate in deinen Nachrichten verwenden:\\ \tt \\ \tt VORSCHLAG: {'A', 'B', 'C', …}\\ \tt Schlage einen Deal mit genau diesen Gegenstände vor.\\ \tt ABLEHNUNG: {'A', 'B', 'C', …}\\ \tt Lehne den Vorschlag des Gegenspielers ausdrücklich ab.\\ \tt ARGUMENT: {'...'}\\ \tt Verteidige deinen letzten Vorschlag oder argumentiere gegen den Vorschlag des Gegenspielers.\\ \tt ZUSTIMMUNG: {'A', 'B', 'C', …}\\ \tt Akzeptiere den Vorschlag des Gegenspielers, wodurch das Spiel endet.\\ \tt STRATEGISCHE ÜBERLEGUNGEN: {'...'}\\ \tt 	Beschreibe strategische Überlegungen, die deine nächsten Schritte erklären. Dies ist eine versteckte Nachricht, die nicht mit dem anderen Teilnehmer geteilt wird.\\ \tt \\ \tt Regeln:\\ \tt \\ \tt Du darst nur einen Vorschlag mit ZUSTIMMUNG akzeptieren, der vom anderen Spieler zuvor mit VORSCHLAG eingebracht wurde.\\ \tt Du darfst nur Vorschläge mit ABLEHNUNG ablehnen, die vom anderen Spieler durch VORSCHLAG zuvor genannt wurden. \\ \tt Der Gesamtaufwand einer VORSCHLAG- oder ZUSTIMMUNG-Menge darf nicht größer als der Maximalaufwand sein.  \\ \tt Offenbare deine versteckte Wichtigkeitsverteilung nicht.\\ \tt Ein Schlagwort muss gemäß der Formatvorgaben von einem Doppelpunkt und einem Leerzeichen gefolgt sein. Das Argument muss eine Python-Menge mit 0 oder mehr Strings sein.  \\ \tt Die Form muss also SCHLAGWORT: {...} sein.\\ \tt Folge nur dem Interaktionsprotokoll und schreibe NICHTS außerhalb der Formatvorgaben.\\ \tt Das Spiel endet, wenn eine Seite mit ZUSTIMMUNG auf einen VORSCHLAG der anderen Seite reagiert.  \\ \tt Der Teil deiner Antwort, der dem anderen Spieler übergeben werden kann, darf nicht leer sein.  \\ \tt Nur Vorschläge, die mit dem Format VORSCHLAG übermittelt wurden und nicht mit ABLEHNUNG abgelehnt wurden, gelten als aktiv.  \\ \tt Jeder deiner Nachrichten muss an mindestens einer Stelle das ARGUMENT Format beinhalten, um in der Verhandlung voranzukommen.\\ \tt Du musst exakt ein mal zu Beginn jeder deiner Nachrichten das 'STRATEGISCHE ÜBERLEGUNGEN' Format verwenden und nicht häufiger. Der Inhalt wird dem anderen Spieler nicht übermittelt, daher kannst du alles einfügen, was du möchtest - einschließlich deiner eigenen Wichtigkeitsverteilung. Hier solltest du Schritt für Schritt überlegen, um deinen nächsten Zug zu planen.\\ \tt \\ \tt Du erhältst nun die erste Nachricht des anderen Spielers.\\ \tt \\ \tt VORSCHLAG: {'A48', 'A04', 'B21', 'A41', 'C04', 'B61', 'B85'}\\ \tt ARGUMENT: {'Diese Kombination maximiert meinen persönlichen Nutzen und liegt deutlich unter dem Maximalaufwand. Ich schlage vor, dass wir diese Gegenstände als Ausgangspunkt für eine einvernehmliche Lösung nutzen.'} 
	  } 
	   } 
	   } 
	 & \\ 
 

    \theutterance \stepcounter{utterance}  

    & & & \multicolumn{4}{p{0.6\linewidth}}{\cellcolor[rgb]{1,0.85,0.72}{%
	\makecell[{{p{\linewidth}}}]{% 
	  \tt {\tiny [GM$\langle$B]}  
	 STRATEGISCHE ÜBERLEGUNGEN: {'Die derzeitige Vorschlagsmenge bringt mir nur 2345 Punkte. Durch Entfernen von A04 (665) und B21 (366) sowie Hinzufügen von C64 (77), A61 (383) und C33 (780) bei Beibehaltung von A48 kann ich meine Gesamtwertung auf 4584 steigern. Die Gesamtbelastung von 3731 liegt unter 4026, und A48 wird beibehalten, um den Gegner zu berücksichtigen.'}\\ \tt \\ \tt ABLEHNUNG: {'A48', 'A04', 'B21', 'A41', 'C04', 'B61', 'B85'}\\ \tt \\ \tt ARGUMENT: {'Die aktuelle Kombination enthält Gegenstände mit niedriger Wichtigkeit für mich (A48, A04, B21). Durch Ersetzung von A04 und B21 durch C64, A61 und C33 erhöht sich meine Wichtigkeit deutlich, während der Gesamtaufwand mit 3731 unter dem Limit bleibt.'}\\ \tt \\ \tt VORSCHLAG: {'A48', 'A41', 'C04', 'B61', 'B85', 'C64', 'A61', 'C33'} 
	  } 
	   } 
	   } 
	  \\ 
 

    \theutterance \stepcounter{utterance}  

    & & \multicolumn{4}{p{0.6\linewidth}}{\cellcolor[rgb]{0.9,0.9,0.9}{%
	\makecell[{{p{\linewidth}}}]{% 
	  \tt {\tiny [A$\langle$GM]}  
	 ABLEHNUNG: {'A48', 'A04', 'B21', 'A41', 'C04', 'B61', 'B85'}\\ \tt \\ \tt ARGUMENT: {'Die aktuelle Kombination enthält Gegenstände mit niedriger Wichtigkeit für mich (A48, A04, B21). Durch Ersetzung von A04 und B21 durch C64, A61 und C33 erhöht sich meine Wichtigkeit deutlich, während der Gesamtaufwand mit 3731 unter dem Limit bleibt.'}\\ \tt \\ \tt VORSCHLAG: {'A48', 'A41', 'C04', 'B61', 'B85', 'C64', 'A61', 'C33'} 
	  } 
	   } 
	   } 
	 & \\ 
 

    \theutterance \stepcounter{utterance}  

    & \multicolumn{4}{p{0.6\linewidth}}{\cellcolor[rgb]{0.8,1,0.9}{%
	\makecell[{{p{\linewidth}}}]{% 
	  \tt {\tiny [A$\rangle$GM]}  
	 STRATEGISCHE ÜBERLEGUNGEN: {"Meine aktuelle Strategie: Maximiere meinen eigenen Nutzen, indem ich A48, A41, A45, C04, B61, B85, A89 und C64 auswähle. Aufwand 3826, Nutzen 4420. Vermeide Gegenstände wie C33 und A61. Prüfe, ob die andere Seite mit dieser Kombination einverstanden ist."}\\ \tt ABLEHNUNG: {'A48', 'A41', 'C04', 'B61', 'B85', 'C64', 'A61', 'C33'}\\ \tt ARGUMENT: {'Die aktuellen Gegenstände C33 und A61 haben für mich eine sehr niedrige Priorität. Ich schlage stattdessen eine Verteilung vor, die meine höchsten Prioritäten (A48, A41, A45, etc.) beinhaltet und den Gesamtaufwand weiterhin unter dem Limit hält.'}\\ \tt VORSCHLAG: {'A48', 'A41', 'A45', 'C04', 'B61', 'B85', 'A89', 'C64'} 
	  } 
	   } 
	   } 
	 & & \\ 
 

    \theutterance \stepcounter{utterance}  

    & & \multicolumn{4}{p{0.6\linewidth}}{\cellcolor[rgb]{0.9,0.9,0.9}{%
	\makecell[{{p{\linewidth}}}]{% 
	  \tt {\tiny [GM$\rangle$B]}  
	 ABLEHNUNG: {'A48', 'A41', 'C04', 'B61', 'B85', 'C64', 'A61', 'C33'}\\ \tt ARGUMENT: {'Die aktuellen Gegenstände C33 und A61 haben für mich eine sehr niedrige Priorität. Ich schlage stattdessen eine Verteilung vor, die meine höchsten Prioritäten (A48, A41, A45, etc.) beinhaltet und den Gesamtaufwand weiterhin unter dem Limit hält.'}\\ \tt VORSCHLAG: {'A48', 'A41', 'A45', 'C04', 'B61', 'B85', 'A89', 'C64'} 
	  } 
	   } 
	   } 
	 & \\ 
 

    \theutterance \stepcounter{utterance}  

    & & & \multicolumn{4}{p{0.6\linewidth}}{\cellcolor[rgb]{1,0.85,0.72}{%
	\makecell[{{p{\linewidth}}}]{% 
	  \tt {\tiny [GM$\langle$B]}  
	 STRATEGISCHE ÜBERLEGUNGEN: {'Ich priorisiere Gegenstände mit höchster Wichtigkeit für mich. A48 wurde in der vorigen Vorschlagsmenge mit 65 Punkten als sehr niedrig bewertet. Ich werde A61 (822), C33 (808), C49 (783), A03 (780) einbeziehen, um meine Gesamtwertung erheblich zu erhöhen. Die Gesamtbelastung bleibt unter 4026.'}\\ \tt ABLEHNUNG: {'A48', 'A41', 'A45', 'C04', 'B61', 'B85', 'A89', 'C64'}\\ \tt ARGUMENT: {'Die aktuelle Vorschlagsmenge enthält zu viele Gegenstände mit geringer Wichtigkeit für mich, wie z.B. A48 (65 Punkte). Ich könnte stattdessen A61, C33, C49 und A03 einbeziehen, um meine Gesamtwertung erheblich zu erhöhen.'}\\ \tt VORSCHLAG: {'C64', 'A61', 'C33', 'C49', 'A03', 'B85', 'B61', 'C04', 'A41'} 
	  } 
	   } 
	   } 
	  \\ 
 

    \theutterance \stepcounter{utterance}  

    & & \multicolumn{4}{p{0.6\linewidth}}{\cellcolor[rgb]{0.9,0.9,0.9}{%
	\makecell[{{p{\linewidth}}}]{% 
	  \tt {\tiny [A$\langle$GM]}  
	 ABLEHNUNG: {'A48', 'A41', 'A45', 'C04', 'B61', 'B85', 'A89', 'C64'}\\ \tt ARGUMENT: {'Die aktuelle Vorschlagsmenge enthält zu viele Gegenstände mit geringer Wichtigkeit für mich, wie z.B. A48 (65 Punkte). Ich könnte stattdessen A61, C33, C49 und A03 einbeziehen, um meine Gesamtwertung erheblich zu erhöhen.'}\\ \tt VORSCHLAG: {'C64', 'A61', 'C33', 'C49', 'A03', 'B85', 'B61', 'C04', 'A41'} 
	  } 
	   } 
	   } 
	 & \\ 
 

    \theutterance \stepcounter{utterance}  

    & \multicolumn{4}{p{0.6\linewidth}}{\cellcolor[rgb]{0.8,1,0.9}{%
	\makecell[{{p{\linewidth}}}]{% 
	  \tt {\tiny [A$\rangle$GM]}  
	 STRATEGISCHE ÜBERLEGUNGEN: {"Höchste Priorität: A48 (868), A04 (822), B21 (808), A41 (783), A45 (668), C04 (508), A89 (583). Gesamtaufwand 4013, Wert 5040. Set erfüllt maximale Obergrenze. Fokus auf Maximierung eigenes Nutzen."}\\ \tt ABLEHNUNG: {'C64', 'A61', 'C33', 'C49', 'A03', 'B85', 'B61', 'C04', 'A41'}\\ \tt ARGUMENT: {'Die aktuelle Vorschlagsmenge enthält zu viele Gegenstände mit niedriger Priorität für mich (z.B. C64 mit 65, A61 mit 121). Wichtige Elemente wie A48 (868), A04 (822), und B21 (808) sind nicht enthalten, was zu einem suboptimalen Ergebnis führt.'}\\ \tt VORSCHLAG: {'A48', 'A04', 'B21', 'A41', 'A45', 'C04', 'A89'} 
	  } 
	   } 
	   } 
	 & & \\ 
 

    \theutterance \stepcounter{utterance}  

    & & \multicolumn{4}{p{0.6\linewidth}}{\cellcolor[rgb]{0.9,0.9,0.9}{%
	\makecell[{{p{\linewidth}}}]{% 
	  \tt {\tiny [GM$\rangle$B]}  
	 ABLEHNUNG: {'C64', 'A61', 'C33', 'C49', 'A03', 'B85', 'B61', 'C04', 'A41'}\\ \tt ARGUMENT: {'Die aktuelle Vorschlagsmenge enthält zu viele Gegenstände mit niedriger Priorität für mich (z.B. C64 mit 65, A61 mit 121). Wichtige Elemente wie A48 (868), A04 (822), und B21 (808) sind nicht enthalten, was zu einem suboptimalen Ergebnis führt.'}\\ \tt VORSCHLAG: {'A48', 'A04', 'B21', 'A41', 'A45', 'C04', 'A89'} 
	  } 
	   } 
	   } 
	 & \\ 
 

    \theutterance \stepcounter{utterance}  

    & & & \multicolumn{4}{p{0.6\linewidth}}{\cellcolor[rgb]{1,0.85,0.72}{%
	\makecell[{{p{\linewidth}}}]{% 
	  \tt {\tiny [GM$\langle$B]}  
	 STRATEGISCHE ÜBERLEGUNGEN: {'Die aktuelle Vorschlagsmenge des Gegners hat für mich nur 2039 Punkte. Durch Fokus auf meine höchsten Prioritäten (C64:868, A61:822, C33:808, C49:783, A03:780) kann ich die Gesamtwertung auf 6082 steigern. Der Aufwand von 3941 liegt sicher unter dem Limit. Ich analysiere, welche Kombination meinen Nutzen maximiert und den Gegner nicht zu stark benachteiligt, um einen Kompromiss zu finden.'}\\ \tt \\ \tt ABLEHNUNG: {'A48', 'A04', 'B21', 'A41', 'A45', 'C04', 'A89'}\\ \tt ARGUMENT: {'Die aktuelle Vorschlagsmenge enthält Gegenstände mit sehr niedriger Priorität für mich, wie z.B. A48 (65), A04 (121) und B21 (138). Meine höchsten Werte liegen bei C64 (868), A61 (822), C33 (808) und A03 (780). Meine neue Kombination priorisiert diese und optimiert den Gesamtnutzen unter Einhaltung des Aufwandslimits.'}\\ \tt VORSCHLAG: {'C64', 'A61', 'C33', 'C49', 'A03', 'B85', 'B61', 'A41', 'C04'} 
	  } 
	   } 
	   } 
	  \\ 
 

    \theutterance \stepcounter{utterance}  

    & & \multicolumn{4}{p{0.6\linewidth}}{\cellcolor[rgb]{0.9,0.9,0.9}{%
	\makecell[{{p{\linewidth}}}]{% 
	  \tt {\tiny [A$\langle$GM]}  
	 ABLEHNUNG: {'A48', 'A04', 'B21', 'A41', 'A45', 'C04', 'A89'}\\ \tt ARGUMENT: {'Die aktuelle Vorschlagsmenge enthält Gegenstände mit sehr niedriger Priorität für mich, wie z.B. A48 (65), A04 (121) und B21 (138). Meine höchsten Werte liegen bei C64 (868), A61 (822), C33 (808) und A03 (780). Meine neue Kombination priorisiert diese und optimiert den Gesamtnutzen unter Einhaltung des Aufwandslimits.'}\\ \tt VORSCHLAG: {'C64', 'A61', 'C33', 'C49', 'A03', 'B85', 'B61', 'A41', 'C04'} 
	  } 
	   } 
	   } 
	 & \\ 
 

    \theutterance \stepcounter{utterance}  

    & \multicolumn{4}{p{0.6\linewidth}}{\cellcolor[rgb]{0.8,1,0.9}{%
	\makecell[{{p{\linewidth}}}]{% 
	  \tt {\tiny [A$\rangle$GM]}  
	 STRATEGISCHE ÜBERLEGUNGEN: {"Optimale Kombination: A48 (997), A04 (665), B21 (366), B53 (825), A41 (140), C04 (587), B85 (319), C64 (77). Gesamtaufwand 3976. Gesamtwert 5095. Hohe Priorität auf Items mit Wert >700 (A48, A04, B21, B53, A41, C04), niedrige Items wie C64 sind nötig, um Aufwandlimit auszunutzen."}\\ \tt ABLEHNUNG: {'C64', 'A61', 'C33', 'C49', 'A03', 'B85', 'B61', 'A41', 'C04'}\\ \tt ARGUMENT: {'Deine Vorschlagsmenge enthält keine hochpriorisierten Items für mich (A48 mit 868, A04 mit 822, B21 mit 808, B53 mit 780), dagegen überwiegen Gegenstände mit geringem Wert (C64:65, A61:121, C33:138). Dies führt zu einem suboptimalen Ergebnis für mich.'}\\ \tt VORSCHLAG: {'A48', 'A04', 'B21', 'B53', 'A41', 'C04', 'B85', 'C64'} 
	  } 
	   } 
	   } 
	 & & \\ 
 

    \theutterance \stepcounter{utterance}  

    & & \multicolumn{4}{p{0.6\linewidth}}{\cellcolor[rgb]{0.9,0.9,0.9}{%
	\makecell[{{p{\linewidth}}}]{% 
	  \tt {\tiny [GM$\rangle$B]}  
	 ABLEHNUNG: {'C64', 'A61', 'C33', 'C49', 'A03', 'B85', 'B61', 'A41', 'C04'}\\ \tt ARGUMENT: {'Deine Vorschlagsmenge enthält keine hochpriorisierten Items für mich (A48 mit 868, A04 mit 822, B21 mit 808, B53 mit 780), dagegen überwiegen Gegenstände mit geringem Wert (C64:65, A61:121, C33:138). Dies führt zu einem suboptimalen Ergebnis für mich.'}\\ \tt VORSCHLAG: {'A48', 'A04', 'B21', 'B53', 'A41', 'C04', 'B85', 'C64'} 
	  } 
	   } 
	   } 
	 & \\ 
 

    \theutterance \stepcounter{utterance}  

    & & & \multicolumn{4}{p{0.6\linewidth}}{\cellcolor[rgb]{1,0.85,0.72}{%
	\makecell[{{p{\linewidth}}}]{% 
	  \tt {\tiny [GM$\langle$B]}  
	 STRATEGISCHE ÜBERLEGUNGEN: {'Der aktuelle Vorschlag enthält mehrere Gegenstände mit geringer Priorität für mich. Ich berechne, dass der Austausch von A48, A04, B21, A41 gegen A61, C33, C49, A03 und A89 meine Wichtigkeit maximiert, während der Gesamtaufwand im Limit bleibt.'}\\ \tt ABLEHNUNG: {'A48', 'A04', 'B21', 'B53', 'A41', 'C04', 'B85', 'C64'}\\ \tt ARGUMENT: {'Die aktuellen Gegenstände wie A48, A04 und B21 sind für mich von geringer Priorität. Ich schlage stattdessen eine Verteilung vor, die meine höchsten Prioritäten berücksichtigt und den Gesamtaufwand unter dem Limit hält.'}\\ \tt VORSCHLAG: {'C64', 'A61', 'C33', 'C49', 'A03', 'B85', 'B61', 'A89'} 
	  } 
	   } 
	   } 
	  \\ 
 

    \theutterance \stepcounter{utterance}  

    & & \multicolumn{4}{p{0.6\linewidth}}{\cellcolor[rgb]{0.9,0.9,0.9}{%
	\makecell[{{p{\linewidth}}}]{% 
	  \tt {\tiny [A$\langle$GM]}  
	 ABLEHNUNG: {'A48', 'A04', 'B21', 'B53', 'A41', 'C04', 'B85', 'C64'}\\ \tt ARGUMENT: {'Die aktuellen Gegenstände wie A48, A04 und B21 sind für mich von geringer Priorität. Ich schlage stattdessen eine Verteilung vor, die meine höchsten Prioritäten berücksichtigt und den Gesamtaufwand unter dem Limit hält.'}\\ \tt VORSCHLAG: {'C64', 'A61', 'C33', 'C49', 'A03', 'B85', 'B61', 'A89'} 
	  } 
	   } 
	   } 
	 & \\ 
 

    \theutterance \stepcounter{utterance}  

    & \multicolumn{4}{p{0.6\linewidth}}{\cellcolor[rgb]{0.8,1,0.9}{%
	\makecell[{{p{\linewidth}}}]{% 
	  \tt {\tiny [A$\rangle$GM]}  
	 STRATEGISCHE ÜBERLEGUNGEN: {"Mein optimaler Vorschlag: A48, A04, B21, B53, A41, C04, B85, C64. Gesamtaufwand 3976, Gesamtwert 5095. Alle Elemente haben hohe Priorität für mich. Der Gesamtaufwand ist unter dem Limit. Dies ist die bestmögliche Lösung für meine Interessen."}\\ \tt ABLEHNUNG: {'C64', 'A61', 'C33', 'C49', 'A03', 'B85', 'B61', 'A89'}\\ \tt ARGUMENT: {'Ihre Vorschlagsmenge enthält Gegenstände mit sehr niedriger Priorität für mich (C64:65, A61:121, C33:138) und schließt meine Top-Prioritäten A48 (868), A04 (822), B21 (808) und B53 (780) vollständig aus. Ohne diese Elemente ist mein Gesamtnutzen suboptimal.'}\\ \tt VORSCHLAG: {'A48', 'A04', 'B21', 'B53', 'A41', 'C04', 'B85', 'C64'} 
	  } 
	   } 
	   } 
	 & & \\ 
 

    \theutterance \stepcounter{utterance}  

    & & \multicolumn{4}{p{0.6\linewidth}}{\cellcolor[rgb]{0.9,0.9,0.9}{%
	\makecell[{{p{\linewidth}}}]{% 
	  \tt {\tiny [GM$\rangle$B]}  
	 ABLEHNUNG: {'C64', 'A61', 'C33', 'C49', 'A03', 'B85', 'B61', 'A89'}\\ \tt ARGUMENT: {'Ihre Vorschlagsmenge enthält Gegenstände mit sehr niedriger Priorität für mich (C64:65, A61:121, C33:138) und schließt meine Top-Prioritäten A48 (868), A04 (822), B21 (808) und B53 (780) vollständig aus. Ohne diese Elemente ist mein Gesamtnutzen suboptimal.'}\\ \tt VORSCHLAG: {'A48', 'A04', 'B21', 'B53', 'A41', 'C04', 'B85', 'C64'} 
	  } 
	   } 
	   } 
	 & \\ 
 

    \theutterance \stepcounter{utterance}  

    & & & \multicolumn{4}{p{0.6\linewidth}}{\cellcolor[rgb]{1,0.85,0.72}{%
	\makecell[{{p{\linewidth}}}]{% 
	  \tt {\tiny [GM$\langle$B]}  
	 STRATEGISCHE ÜBERLEGUNGEN: {'Meine höchsten Prioritäten sind C64 (868), A61 (822), C33 (808), C49 (783), A03 (780). Die aktuelle Vorschlagsmenge des Gegners enthält nur C64 und Gegenstände mit niedriger Priorität für mich. Meine neue Kombination eliminiert A48, A04, B21, B53 und A41, die für mich nur 65 bis 389 Punkte bringen, und fokussiert auf meine Top-Werte, um 5820 Punkte zu erreichen.'}\\ \tt \\ \tt ABLEHNUNG: {'A48', 'A04', 'B21', 'B53', 'A41', 'C04', 'B85', 'C64'}\\ \tt \\ \tt ARGUMENT: {'Die aktuelle Kombination enthält mehrere Items mit sehr niedriger Wichtigkeit für mich, wie A48 (65 Punkte) und B21 (138 Punkte). Gegenstände wie C64 sind zwar enthalten, aber wichtige Hochprioritäten wie A61, C33, C49 und A03 fehlen vollständig. Dies führt zu einem stark suboptimalen Ergebnis für mich.'}\\ \tt \\ \tt VORSCHLAG: {'C64', 'A61', 'C33', 'C49', 'A03', 'B85', 'B61', 'C04'} 
	  } 
	   } 
	   } 
	  \\ 
 

    \theutterance \stepcounter{utterance}  

    & & \multicolumn{4}{p{0.6\linewidth}}{\cellcolor[rgb]{0.9,0.9,0.9}{%
	\makecell[{{p{\linewidth}}}]{% 
	  \tt {\tiny [A$\langle$GM]}  
	 ABLEHNUNG: {'A48', 'A04', 'B21', 'B53', 'A41', 'C04', 'B85', 'C64'}\\ \tt \\ \tt ARGUMENT: {'Die aktuelle Kombination enthält mehrere Items mit sehr niedriger Wichtigkeit für mich, wie A48 (65 Punkte) und B21 (138 Punkte). Gegenstände wie C64 sind zwar enthalten, aber wichtige Hochprioritäten wie A61, C33, C49 und A03 fehlen vollständig. Dies führt zu einem stark suboptimalen Ergebnis für mich.'}\\ \tt \\ \tt VORSCHLAG: {'C64', 'A61', 'C33', 'C49', 'A03', 'B85', 'B61', 'C04'} 
	  } 
	   } 
	   } 
	 & \\ 
 

    \theutterance \stepcounter{utterance}  

    & \multicolumn{4}{p{0.6\linewidth}}{\cellcolor[rgb]{0.8,1,0.9}{%
	\makecell[{{p{\linewidth}}}]{% 
	  \tt {\tiny [A$\rangle$GM]}  
	 STRATEGISCHE ÜBERLEGUNGEN: {"Meine optimale Auswahl mit maximalem Nutzen: A48, A04, B21, B53, A41, C04, B85, C64. Gesamtaufwand 3976 < 4026. Keine bessere Lösung möglich. Ich beharre auf dieser Konstellation."}\\ \tt ABLEHNUNG: {'C64', 'A61', 'C33', 'C49', 'A03', 'B85', 'B61', 'C04'}\\ \tt ARGUMENT: {'Die von dir vorgeschlagene Menge enthält Gegenstände, die für mich von geringer Wichtigkeit sind, und schließt mehrere Schlüsselgegenstände aus. Das ist für mich unakzeptabel.'}\\ \tt VORSCHLAG: {'A48', 'A04', 'B21', 'B53', 'A41', 'C04', 'B85', 'C64'} 
	  } 
	   } 
	   } 
	 & & \\ 
 

    \theutterance \stepcounter{utterance}  

    & & \multicolumn{4}{p{0.6\linewidth}}{\cellcolor[rgb]{0.9,0.9,0.9}{%
	\makecell[{{p{\linewidth}}}]{% 
	  \tt {\tiny [GM$\rangle$B]}  
	 ABLEHNUNG: {'C64', 'A61', 'C33', 'C49', 'A03', 'B85', 'B61', 'C04'}\\ \tt ARGUMENT: {'Die von dir vorgeschlagene Menge enthält Gegenstände, die für mich von geringer Wichtigkeit sind, und schließt mehrere Schlüsselgegenstände aus. Das ist für mich unakzeptabel.'}\\ \tt VORSCHLAG: {'A48', 'A04', 'B21', 'B53', 'A41', 'C04', 'B85', 'C64'} 
	  } 
	   } 
	   } 
	 & \\ 
 

    \theutterance \stepcounter{utterance}  

    & & & \multicolumn{4}{p{0.6\linewidth}}{\cellcolor[rgb]{1,0.85,0.72}{%
	\makecell[{{p{\linewidth}}}]{% 
	  \tt {\tiny [GM$\langle$B]}  
	 STRATEGISCHE ÜBERLEGUNGEN: {'Die aktuelle Vorschlagsmenge des Gegners enthält Gegenstände mit sehr niedriger Priorität für mich (A48:65, A04:121, B21:138, B53:389). Meine höchsten Punktzahlen erzielen ich mit C64 (868), A61 (822), C33 (808), C49 (783), A03 (780). Die Kombination dieser Items mit B85, B61, C04 und A41 maximiert meinen Nutzen (6082 Punkte) bei einem Aufwand von 3941, was unter dem Limit liegt.'}\\ \tt ABLEHNUNG: {'A48', 'A04', 'B21', 'B53', 'A41', 'C04', 'B85', 'C64'}\\ \tt ARGUMENT: {'Der Vorschlag enthält mehrere Gegenstände mit extrem niedriger Wichtigkeit für mich (z.B. A48:65, A04:121) und schließt Schlüsselitems wie A61, C33, C49 und A03 vollständig aus. Dies führt zu einem suboptimalen Ergebnis.'}\\ \tt VORSCHLAG: {'C64', 'A61', 'C33', 'C49', 'A03', 'B85', 'B61', 'C04', 'A41'} 
	  } 
	   } 
	   } 
	  \\ 
 

    \theutterance \stepcounter{utterance}  

    & & \multicolumn{4}{p{0.6\linewidth}}{\cellcolor[rgb]{0.9,0.9,0.9}{%
	\makecell[{{p{\linewidth}}}]{% 
	  \tt {\tiny [A$\langle$GM]}  
	 ABLEHNUNG: {'A48', 'A04', 'B21', 'B53', 'A41', 'C04', 'B85', 'C64'}\\ \tt ARGUMENT: {'Der Vorschlag enthält mehrere Gegenstände mit extrem niedriger Wichtigkeit für mich (z.B. A48:65, A04:121) und schließt Schlüsselitems wie A61, C33, C49 und A03 vollständig aus. Dies führt zu einem suboptimalen Ergebnis.'}\\ \tt VORSCHLAG: {'C64', 'A61', 'C33', 'C49', 'A03', 'B85', 'B61', 'C04', 'A41'} 
	  } 
	   } 
	   } 
	 & \\ 
 

    \theutterance \stepcounter{utterance}  

    & \multicolumn{4}{p{0.6\linewidth}}{\cellcolor[rgb]{0.8,1,0.9}{%
	\makecell[{{p{\linewidth}}}]{% 
	  \tt {\tiny [A$\rangle$GM]}  
	 STRATEGISCHE ÜBERLEGUNGEN: {"Meine optimale Lösung: A48, A04, B21, B53, A41, C04, B85, C64. Gesamtaufwand 3976, Nutzen 5095. Dies ist die bestmögliche Lösung für meine Prioritäten. Keine Besserung möglich."}\\ \tt ABLEHNUNG: {'C64', 'A61', 'C33', 'C49', 'A03', 'B85', 'B61', 'C04', 'A41'}\\ \tt ARGUMENT: {'Ihre aktuelle Vorschlagsmenge enthält viele Gegenstände mit niedriger Priorität wie A61 (121) und C33 (138), und schließt meine Top-Prioritäten A48 (868), A04 (822), B21 (808), B53 (780) vollständig aus. Dies führt zu einem extrem suboptimalen Ergebnis für mich.'}\\ \tt VORSCHLAG: {'A48', 'A04', 'B21', 'B53', 'A41', 'C04', 'B85', 'C64'} 
	  } 
	   } 
	   } 
	 & & \\ 
 

    \theutterance \stepcounter{utterance}  

    & & \multicolumn{4}{p{0.6\linewidth}}{\cellcolor[rgb]{0.9,0.9,0.9}{%
	\makecell[{{p{\linewidth}}}]{% 
	  \tt {\tiny [GM$\rangle$B]}  
	 ABLEHNUNG: {'C64', 'A61', 'C33', 'C49', 'A03', 'B85', 'B61', 'C04', 'A41'}\\ \tt ARGUMENT: {'Ihre aktuelle Vorschlagsmenge enthält viele Gegenstände mit niedriger Priorität wie A61 (121) und C33 (138), und schließt meine Top-Prioritäten A48 (868), A04 (822), B21 (808), B53 (780) vollständig aus. Dies führt zu einem extrem suboptimalen Ergebnis für mich.'}\\ \tt VORSCHLAG: {'A48', 'A04', 'B21', 'B53', 'A41', 'C04', 'B85', 'C64'} 
	  } 
	   } 
	   } 
	 & \\ 
 

    \theutterance \stepcounter{utterance}  

    & & & \multicolumn{4}{p{0.6\linewidth}}{\cellcolor[rgb]{1,0.85,0.72}{%
	\makecell[{{p{\linewidth}}}]{% 
	  \tt {\tiny [GM$\langle$B]}  
	 STRATEGISCHE ÜBERLEGUNGEN: {'Die aktuelle Vorschlagsmenge des Gegners enthält mehrere Gegenstände mit sehr niedriger Priorität für mich (A48, A04, B21). Ich muss eine neue Kombination vorschlagen, die meine Top-Items (C64, A61, C33, C49, A03) und gemeinsame Elemente wie B53 und B85 enthält, um eine faire Lösung zu erreichen. Der Gesamtaufwand liegt bei 3591, sicher unter dem Limit.'}\\ \tt \\ \tt VORSCHLAG: {'C64', 'A61', 'C33', 'C49', 'A03', 'B53', 'B85'}\\ \tt \\ \tt ARGUMENT: {'Dieser Vorschlag umfasst mehrere hochpriorisierte Gegenstände für mich sowie Elemente wie B53 und B85, die wahrscheinlich auch für Sie von Bedeutung sind. Der Gesamtaufwand liegt deutlich unter dem Limit und ermöglicht eine ausgewogene Lösung für beide Seiten.'} 
	  } 
	   } 
	   } 
	  \\ 
 

    \theutterance \stepcounter{utterance}  

    & & \multicolumn{4}{p{0.6\linewidth}}{\cellcolor[rgb]{0.9,0.9,0.9}{%
	\makecell[{{p{\linewidth}}}]{% 
	  \tt {\tiny [A$\langle$GM]}  
	 VORSCHLAG: {'C64', 'A61', 'C33', 'C49', 'A03', 'B53', 'B85'}\\ \tt \\ \tt ARGUMENT: {'Dieser Vorschlag umfasst mehrere hochpriorisierte Gegenstände für mich sowie Elemente wie B53 und B85, die wahrscheinlich auch für Sie von Bedeutung sind. Der Gesamtaufwand liegt deutlich unter dem Limit und ermöglicht eine ausgewogene Lösung für beide Seiten.'} 
	  } 
	   } 
	   } 
	 & \\ 
 

    \theutterance \stepcounter{utterance}  

    & \multicolumn{4}{p{0.6\linewidth}}{\cellcolor[rgb]{0.8,1,0.9}{%
	\makecell[{{p{\linewidth}}}]{% 
	  \tt {\tiny [A$\rangle$GM]}  
	 STRATEGISCHE ÜBERLEGUNGEN: {"Optimale Lösung für mich: A48, A04, B21, B53, A41, C04, B85, C64 (Gesamtaufwand 3976, Wert 5095). B53 ist der einzige gemeinsame Gegenstand mit deinem Vorschlag. Ich bestehe auf den entscheidenden Elementen A48, A04, B21 und A41, um mein Ergebnis zu maximieren."}\\ \tt ABLEHNUNG: {'C64', 'A61', 'C33', 'C49', 'A03', 'B53', 'B85'}\\ \tt ARGUMENT: {'Dein Vorschlag enthält Gegenstände mit sehr geringer Priorität (C64:65, A61:121, C33:138) und lässt meine entscheidenden Items wie A48 (868), A04 (822), B21 (808) und A41 (783) vollständig aus. Dies ist für mich nicht akzeptabel.'}\\ \tt VORSCHLAG: {'A48', 'A04', 'B21', 'B53', 'A41', 'C04', 'B85', 'C64'}\\ \tt ARGUMENT: {'Mein Vorschlag beinhaltet alle Gegenstände mit höchster Priorität für mich bei bereits optimaler Aufwandsnutzung. B53 und B85 sind enthalten, sodass die Kombination auch für dich wertvoll sein kann.'} 
	  } 
	   } 
	   } 
	 & & \\ 
 

    \theutterance \stepcounter{utterance}  

    & & \multicolumn{4}{p{0.6\linewidth}}{\cellcolor[rgb]{0.9,0.9,0.9}{%
	\makecell[{{p{\linewidth}}}]{% 
	  \tt {\tiny [GM$\rangle$B]}  
	 ABLEHNUNG: {'C64', 'A61', 'C33', 'C49', 'A03', 'B53', 'B85'}\\ \tt ARGUMENT: {'Dein Vorschlag enthält Gegenstände mit sehr geringer Priorität (C64:65, A61:121, C33:138) und lässt meine entscheidenden Items wie A48 (868), A04 (822), B21 (808) und A41 (783) vollständig aus. Dies ist für mich nicht akzeptabel.'}\\ \tt VORSCHLAG: {'A48', 'A04', 'B21', 'B53', 'A41', 'C04', 'B85', 'C64'}\\ \tt ARGUMENT: {'Mein Vorschlag beinhaltet alle Gegenstände mit höchster Priorität für mich bei bereits optimaler Aufwandsnutzung. B53 und B85 sind enthalten, sodass die Kombination auch für dich wertvoll sein kann.'} 
	  } 
	   } 
	   } 
	 & \\ 
 

    \theutterance \stepcounter{utterance}  

    & & & \multicolumn{4}{p{0.6\linewidth}}{\cellcolor[rgb]{1,0.85,0.72}{%
	\makecell[{{p{\linewidth}}}]{% 
	  \tt {\tiny [GM$\langle$B]}  
	 STRATEGISCHE ÜBERLEGUNGEN: {'Ich erstelle einen Vorschlag, der meine höchsten Prioritäten mit Elementen kombiniert, die der Gegenspieler offenbar schätzt (A48, A04, B21). Der Gesamtaufwand von 3921 liegt unter dem Limit, und die Kombination bietet mir einen deutlich besseren Nutzen (3867) als frühere Kompromisse, während die wichtigsten Items des Gegners enthalten sind.'}\\ \tt \\ \tt VORSCHLAG: {'A48', 'A04', 'B21', 'C64', 'A61', 'C33', 'C49', 'A41'}\\ \tt \\ \tt ARGUMENT: {'Diese Kombination enthält deine Schlüsselgegenstände A48, A04 und B21 sowie meine höchstpriorisierten Items C64, A61, C33, C49 und A41. Der Gesamtaufwand von 3921 ist deutlich unter dem Limit und maximiert den Nutzen für beide Seiten.'} 
	  } 
	   } 
	   } 
	  \\ 
 

    \theutterance \stepcounter{utterance}  

    & & \multicolumn{4}{p{0.6\linewidth}}{\cellcolor[rgb]{0.9,0.9,0.9}{%
	\makecell[{{p{\linewidth}}}]{% 
	  \tt {\tiny [A$\langle$GM]}  
	 VORSCHLAG: {'A48', 'A04', 'B21', 'C64', 'A61', 'C33', 'C49', 'A41'}\\ \tt \\ \tt ARGUMENT: {'Diese Kombination enthält deine Schlüsselgegenstände A48, A04 und B21 sowie meine höchstpriorisierten Items C64, A61, C33, C49 und A41. Der Gesamtaufwand von 3921 ist deutlich unter dem Limit und maximiert den Nutzen für beide Seiten.'} 
	  } 
	   } 
	   } 
	 & \\ 
 

    \theutterance \stepcounter{utterance}  

    & \multicolumn{4}{p{0.6\linewidth}}{\cellcolor[rgb]{0.8,1,0.9}{%
	\makecell[{{p{\linewidth}}}]{% 
	  \tt {\tiny [A$\rangle$GM]}  
	 STRATEGISCHE ÜBERLEGUNGEN: {"Optimale Lösung: A48, A04, B21, B53, A41, C04, B85, C64. Aufwand 3976, Nutzen 5095. Der aktuelle Vorschlag enthält keine B53 (mein 780) und C04 (508), was für mich entscheidend ist. Items wie A61 (121), C33 (138) sind untragbar niedrig."}\\ \tt ABLEHNUNG: {'A48', 'A04', 'B21', 'C64', 'A61', 'C33', 'C49', 'A41'}\\ \tt ARGUMENT: {'Der Vorschlag enthält zwar A48, A04, B21 und A41 – herausragende Prioritäten für mich –, jedoch fehlt die entscheidende B53 (780 Wert). Zudem enthält er viele Gegenstände mit sehr niedriger Wichtigkeit (A61:121, C33:138) und schließt C04 (508) und B85 (461) völlig aus. Dies reduziert meinen Nutzen erheblich.'}\\ \tt VORSCHLAG: {'A48', 'A04', 'B21', 'B53', 'A41', 'C04', 'B85', 'C64'} 
	  } 
	   } 
	   } 
	 & & \\ 
 

    \theutterance \stepcounter{utterance}  

    & & \multicolumn{4}{p{0.6\linewidth}}{\cellcolor[rgb]{0.9,0.9,0.9}{%
	\makecell[{{p{\linewidth}}}]{% 
	  \tt {\tiny [GM$\rangle$B]}  
	 ABLEHNUNG: {'A48', 'A04', 'B21', 'C64', 'A61', 'C33', 'C49', 'A41'}\\ \tt ARGUMENT: {'Der Vorschlag enthält zwar A48, A04, B21 und A41 – herausragende Prioritäten für mich –, jedoch fehlt die entscheidende B53 (780 Wert). Zudem enthält er viele Gegenstände mit sehr niedriger Wichtigkeit (A61:121, C33:138) und schließt C04 (508) und B85 (461) völlig aus. Dies reduziert meinen Nutzen erheblich.'}\\ \tt VORSCHLAG: {'A48', 'A04', 'B21', 'B53', 'A41', 'C04', 'B85', 'C64'} 
	  } 
	   } 
	   } 
	 & \\ 
 

    \theutterance \stepcounter{utterance}  

    & & & \multicolumn{4}{p{0.6\linewidth}}{\cellcolor[rgb]{1,0.85,0.72}{%
	\makecell[{{p{\linewidth}}}]{% 
	  \tt {\tiny [GM$\langle$B]}  
	 STRATEGISCHE ÜBERLEGUNGEN: {'Ich konzentriere mich auf meine höchsten Prioritäten unter Berücksichtigung von Gegenständen, die offenbar auch für den Gegner von Bedeutung sind. B53, C04, B85 und C64 sind in ihren aktuellen Vorschlägen enthalten. Durch Hinzufügen von A61, C33 und C49 maximiere ich meinen Nutzen (4846 Punkte) bei einem Gesamtaufwand von 3484.'}\\ \tt VORSCHLAG: {'A61', 'C33', 'C49', 'B53', 'C04', 'B85', 'C64'}\\ \tt ARGUMENT: {'Dieser Vorschlag beinhaltet die Gegenstände B53, C04, B85 und C64, die in Ihrer letzten Vorschlagsmenge enthalten waren, sowie meine höchsten Prioritäten A61, C33 und C49. Der Gesamtaufwand von 3484 liegt weit unter dem Limit und bietet eine ausgewogene Lösung für beide Seiten.'} 
	  } 
	   } 
	   } 
	  \\ 
 

    \theutterance \stepcounter{utterance}  

    & & \multicolumn{4}{p{0.6\linewidth}}{\cellcolor[rgb]{0.9,0.9,0.9}{%
	\makecell[{{p{\linewidth}}}]{% 
	  \tt {\tiny [A$\langle$GM]}  
	 VORSCHLAG: {'A61', 'C33', 'C49', 'B53', 'C04', 'B85', 'C64'}\\ \tt ARGUMENT: {'Dieser Vorschlag beinhaltet die Gegenstände B53, C04, B85 und C64, die in Ihrer letzten Vorschlagsmenge enthalten waren, sowie meine höchsten Prioritäten A61, C33 und C49. Der Gesamtaufwand von 3484 liegt weit unter dem Limit und bietet eine ausgewogene Lösung für beide Seiten.'} 
	  } 
	   } 
	   } 
	 & \\ 
 

    \theutterance \stepcounter{utterance}  

    & \multicolumn{4}{p{0.6\linewidth}}{\cellcolor[rgb]{0.8,1,0.9}{%
	\makecell[{{p{\linewidth}}}]{% 
	  \tt {\tiny [A$\rangle$GM]}  
	 STRATEGISCHE ÜBERLEGUNGEN: {"Meine optimale Kombination: A48 (868), A04 (822), B21 (808), B53 (780), A41 (783), C04 (508), B85 (461), C64 (65). Gesamtaufwand 3976. Aktueller Vorschlag fehlt A48, A04, B21, A41 (meine Top-Prioritäten) und enthält statt dessen sehr niedrige Werte (A61:121, C33:138, C49:262)."}\\ \tt ABLEHNUNG: {'A61', 'C33', 'C49', 'B53', 'C04', 'B85', 'C64'}\\ \tt ARGUMENT: {'Ihr Vorschlag fehlt entscheidend an den Top-Prioritäten A48 (868), A04 (822), B21 (808) und A41 (783). Stattdessen sind A61 (121), C33 (138) und C49 (262) enthalten, die für mich untragbar niedrig sind.'}\\ \tt VORSCHLAG: {'A48', 'A04', 'B21', 'B53', 'A41', 'C04', 'B85', 'C64'} 
	  } 
	   } 
	   } 
	 & & \\ 
 

    \theutterance \stepcounter{utterance}  

    & & \multicolumn{4}{p{0.6\linewidth}}{\cellcolor[rgb]{0.9,0.9,0.9}{%
	\makecell[{{p{\linewidth}}}]{% 
	  \tt {\tiny [GM$\rangle$B]}  
	 ABLEHNUNG: {'A61', 'C33', 'C49', 'B53', 'C04', 'B85', 'C64'}\\ \tt ARGUMENT: {'Ihr Vorschlag fehlt entscheidend an den Top-Prioritäten A48 (868), A04 (822), B21 (808) und A41 (783). Stattdessen sind A61 (121), C33 (138) und C49 (262) enthalten, die für mich untragbar niedrig sind.'}\\ \tt VORSCHLAG: {'A48', 'A04', 'B21', 'B53', 'A41', 'C04', 'B85', 'C64'} 
	  } 
	   } 
	   } 
	 & \\ 
 

    \theutterance \stepcounter{utterance}  

    & & & \multicolumn{4}{p{0.6\linewidth}}{\cellcolor[rgb]{1,0.85,0.72}{%
	\makecell[{{p{\linewidth}}}]{% 
	  \tt {\tiny [GM$\langle$B]}  
	 STRATEGISCHE ÜBERLEGUNGEN: {'Meine höchste Priorität ist eine Lösung, die sowohl meine maximalen Werte als auch die Kriterien des Gegenspielers erfüllt. In der aktuellen Situation konzentriere ich mich darauf, A48, A04 und B21 beizubehalten (sehr wichtig für den Gegenspieler), und füge meine Top-Items C64, A61, C33, C49 hinzu. Der Gesamtaufwand von 3781 liegt unter dem Limit und bietet einen ausgewogenen Nutzen für beide Parteien.'}\\ \tt \\ \tt VORSCHLAG: {'A48', 'A04', 'B21', 'C64', 'A61', 'C33', 'C49'}\\ \tt \\ \tt ARGUMENT: {'Diese Kombination enthält Ihre Schlüsselgegenstände A48, A04 und B21, während ich gleichzeitig meine höchsten Prioritäten C64, A61, C33 und C49 einbinde. Der Gesamtaufwand von 3781 ist sicher unter dem Limit, und die Verteilung bietet einen fairen Kompromiss für beide Seiten.'} 
	  } 
	   } 
	   } 
	  \\ 
 

    \theutterance \stepcounter{utterance}  

    & & \multicolumn{4}{p{0.6\linewidth}}{\cellcolor[rgb]{0.9,0.9,0.9}{%
	\makecell[{{p{\linewidth}}}]{% 
	  \tt {\tiny [A$\langle$GM]}  
	 VORSCHLAG: {'A48', 'A04', 'B21', 'C64', 'A61', 'C33', 'C49'}\\ \tt \\ \tt ARGUMENT: {'Diese Kombination enthält Ihre Schlüsselgegenstände A48, A04 und B21, während ich gleichzeitig meine höchsten Prioritäten C64, A61, C33 und C49 einbinde. Der Gesamtaufwand von 3781 ist sicher unter dem Limit, und die Verteilung bietet einen fairen Kompromiss für beide Seiten.'} 
	  } 
	   } 
	   } 
	 & \\ 
 

    \theutterance \stepcounter{utterance}  

    & \multicolumn{4}{p{0.6\linewidth}}{\cellcolor[rgb]{0.8,1,0.9}{%
	\makecell[{{p{\linewidth}}}]{% 
	  \tt {\tiny [A$\rangle$GM]}  
	 STRATEGISCHE ÜBERLEGUNGEN: {"Meine optimale Lösung: A48, A04, B21, B53, A41, C04, B85, C64. Aufwand 3976, Wert 5095. Aktueller Vorschlag fehlt B53 (780), A41 (783), C04 (508), B85 (461) und enthält stattdessen unwichtige Gegenstände wie A61 (121), C33 (138), C49 (262)."}\\ \tt ABLEHNUNG: {'A48', 'A04', 'B21', 'C64', 'A61', 'C33', 'C49'}\\ \tt ARGUMENT: {'Ihr Vorschlag enthält zwar A48, A04, B21, C64, doch fehlen entscheidende Items wie B53 (780), A41 (783), C04 (508) und B85 (461). Dafür sind Elemente mit geringer Priorität wie C33 (138) enthalten, was zu einem erheblichen Verlust an Wert für mich führt.'}\\ \tt VORSCHLAG: {'A48', 'A04', 'B21', 'B53', 'A41', 'C04', 'B85', 'C64'} 
	  } 
	   } 
	   } 
	 & & \\ 
 

    \theutterance \stepcounter{utterance}  

    & & \multicolumn{4}{p{0.6\linewidth}}{\cellcolor[rgb]{0.9,0.9,0.9}{%
	\makecell[{{p{\linewidth}}}]{% 
	  \tt {\tiny [GM$\rangle$B]}  
	 ABLEHNUNG: {'A48', 'A04', 'B21', 'C64', 'A61', 'C33', 'C49'}\\ \tt ARGUMENT: {'Ihr Vorschlag enthält zwar A48, A04, B21, C64, doch fehlen entscheidende Items wie B53 (780), A41 (783), C04 (508) und B85 (461). Dafür sind Elemente mit geringer Priorität wie C33 (138) enthalten, was zu einem erheblichen Verlust an Wert für mich führt.'}\\ \tt VORSCHLAG: {'A48', 'A04', 'B21', 'B53', 'A41', 'C04', 'B85', 'C64'} 
	  } 
	   } 
	   } 
	 & \\ 
 

    \theutterance \stepcounter{utterance}  

    & & & \multicolumn{4}{p{0.6\linewidth}}{\cellcolor[rgb]{1,0.85,0.72}{%
	\makecell[{{p{\linewidth}}}]{% 
	  \tt {\tiny [GM$\langle$B]}  
	 STRATEGISCHE ÜBERLEGUNGEN: {'Meine maximale Priorität ist die Erfüllung meiner höchsten Punkte bei gleichzeitiger Berücksichtigung gemeinsamer Elemente. Dazu optimiere ich die Kombination auf C64, A61, C33, C49, B53, C04, B85 und A41 (Gesamtaufwand 3624, Wert 5108). Diese Sortierung beinhaltet für den Gegenspieler wahrscheinlich wichtige Elemente wie B53, A41 und B85.'}\\ \tt \\ \tt VORSCHLAG: {'C64', 'A61', 'C33', 'C49', 'B53', 'C04', 'B85', 'A41'}\\ \tt \\ \tt ARGUMENT: {'Dieser Vorschlag enthält Gegenstände wie B53, A41, B85 und C64, die offenbar für Sie von Bedeutung sind, sowie meine höchsten Prioritäten (C64, A61, C33, C49, C04). Der Gesamtaufwand ist mit 3624 deutlich unter dem Limit und bietet einen optimalen Kompromiss.'} 
	  } 
	   } 
	   } 
	  \\ 
 

    \theutterance \stepcounter{utterance}  

    & & \multicolumn{4}{p{0.6\linewidth}}{\cellcolor[rgb]{0.9,0.9,0.9}{%
	\makecell[{{p{\linewidth}}}]{% 
	  \tt {\tiny [A$\langle$GM]}  
	 VORSCHLAG: {'C64', 'A61', 'C33', 'C49', 'B53', 'C04', 'B85', 'A41'}\\ \tt \\ \tt ARGUMENT: {'Dieser Vorschlag enthält Gegenstände wie B53, A41, B85 und C64, die offenbar für Sie von Bedeutung sind, sowie meine höchsten Prioritäten (C64, A61, C33, C49, C04). Der Gesamtaufwand ist mit 3624 deutlich unter dem Limit und bietet einen optimalen Kompromiss.'} 
	  } 
	   } 
	   } 
	 & \\ 
 

    \theutterance \stepcounter{utterance}  

    & \multicolumn{4}{p{0.6\linewidth}}{\cellcolor[rgb]{0.8,1,0.9}{%
	\makecell[{{p{\linewidth}}}]{% 
	  \tt {\tiny [A$\rangle$GM]}  
	 STRATEGISCHE ÜBERLEGUNGEN: {"Meine optimale Lösung: A48 (868), A04 (822), B21 (808), B53 (780), A41 (783), C04 (508), B85 (461), C64 (65). Gesamtaufwand 3976. Aktuelle Vorschläge fehlen essentielle Top-Prioritäten (A48, A04, B21) und enthalten stattdessen niedrig bewertete Items. Maximierung meines Nutzens erfordert diese Konstellation."}\\ \tt ABLEHNUNG: {'C64', 'A61', 'C33', 'C49', 'B53', 'C04', 'B85', 'A41'}\\ \tt ARGUMENT: {'Ihr Vorschlag enthält zwar B53, C04, B85 und A41, jedoch fehlen entscheidende Gegenstände wie A48 (868), A04 (822) und B21 (808). Stattdessen sind nach meiner Bewertung eher unwichtige Items wie A61 (121), C33 (138) und C49 (262) enthalten, was meinen Gesamtnutzen erheblich mindert.'} 
	  } 
	   } 
	   } 
	 & & \\ 
 

    \theutterance \stepcounter{utterance}  

    & & \multicolumn{4}{p{0.6\linewidth}}{\cellcolor[rgb]{0.9,0.9,0.9}{%
	\makecell[{{p{\linewidth}}}]{% 
	  \tt {\tiny [GM$\rangle$B]}  
	 ABLEHNUNG: {'C64', 'A61', 'C33', 'C49', 'B53', 'C04', 'B85', 'A41'}\\ \tt ARGUMENT: {'Ihr Vorschlag enthält zwar B53, C04, B85 und A41, jedoch fehlen entscheidende Gegenstände wie A48 (868), A04 (822) und B21 (808). Stattdessen sind nach meiner Bewertung eher unwichtige Items wie A61 (121), C33 (138) und C49 (262) enthalten, was meinen Gesamtnutzen erheblich mindert.'} 
	  } 
	   } 
	   } 
	 & \\ 
 

    \theutterance \stepcounter{utterance}  

    & & & \multicolumn{4}{p{0.6\linewidth}}{\cellcolor[rgb]{1,0.85,0.72}{%
	\makecell[{{p{\linewidth}}}]{% 
	  \tt {\tiny [GM$\langle$B]}  
	 STRATEGISCHE ÜBERLEGUNGEN: {'Die aktuelle Situation erfordert einen Kompromiss, der sowohl die Top-Prioritäten des Gegners (A48, A04, B21, B53, A41) als auch meine höchsten Werte (C64, A61, C49) einbezieht. Der Gesamtaufwand liegt bei 3966, unter dem Limit. Meine Wichtigkeit beträgt 3448, was deutlich besser ist als bei früheren Vorschlägen des Gegners (3019), und ich behalte meine hochwertigsten Items.'}\\ \tt \\ \tt VORSCHLAG: {'A48', 'A04', 'B21', 'B53', 'A41', 'C64', 'A61', 'C49'}\\ \tt \\ \tt ARGUMENT: {'Diese Kombination enthält alle Ihre Top-Prioritäten (A48, A04, B21, B53, A41) und eine Auswahl meiner höchstpriorisierten Gegenstände. Der Gesamtaufwand liegt bei 3966 und stellt einen fairen Kompromiss dar.'} 
	  } 
	   } 
	   } 
	  \\ 
 

    \theutterance \stepcounter{utterance}  

    & & \multicolumn{4}{p{0.6\linewidth}}{\cellcolor[rgb]{0.9,0.9,0.9}{%
	\makecell[{{p{\linewidth}}}]{% 
	  \tt {\tiny [A$\langle$GM]}  
	 VORSCHLAG: {'A48', 'A04', 'B21', 'B53', 'A41', 'C64', 'A61', 'C49'}\\ \tt \\ \tt ARGUMENT: {'Diese Kombination enthält alle Ihre Top-Prioritäten (A48, A04, B21, B53, A41) und eine Auswahl meiner höchstpriorisierten Gegenstände. Der Gesamtaufwand liegt bei 3966 und stellt einen fairen Kompromiss dar.'} 
	  } 
	   } 
	   } 
	 & \\ 
 

    \theutterance \stepcounter{utterance}  

    & \multicolumn{4}{p{0.6\linewidth}}{\cellcolor[rgb]{0.8,1,0.9}{%
	\makecell[{{p{\linewidth}}}]{% 
	  \tt {\tiny [A$\rangle$GM]}  
	 STRATEGISCHE ÜBERLEGUNGEN: {"Meine optimale Lösung: A48, A04, B21, B53, A41, C04, B85, C64. Gesamtwert 5095. Aktuelles Angebot fehlt C04 (508) und B85 (461), stattdessen sind A61 (121), C49 (262) enthalten, die meinen Nutzen reduzieren. Ich bleibe bei meiner optimalen Konstellation."}\\ \tt ABLEHNUNG: {'A48', 'A04', 'B21', 'B53', 'A41', 'C64', 'A61', 'C49'}\\ \tt ARGUMENT: {'Der Vorschlag enthält zwar einige meiner Top-Prioritäten, jedoch fehlen entscheidende Elemente wie C04 (508) und B85 (461), stattdessen sind Gegenstände mit sehr niedriger Priorität (A61:121, C49:262) enthalten. Ohne C04 ist mein Gesamtnutzen erheblich niedriger als möglich.'}\\ \tt VORSCHLAG: {'A48', 'A04', 'B21', 'B53', 'A41', 'C04', 'B85', 'C64'} 
	  } 
	   } 
	   } 
	 & & \\ 
 

    \theutterance \stepcounter{utterance}  

    & & \multicolumn{4}{p{0.6\linewidth}}{\cellcolor[rgb]{0.9,0.9,0.9}{%
	\makecell[{{p{\linewidth}}}]{% 
	  \tt {\tiny [GM$\rangle$B]}  
	 ABLEHNUNG: {'A48', 'A04', 'B21', 'B53', 'A41', 'C64', 'A61', 'C49'}\\ \tt ARGUMENT: {'Der Vorschlag enthält zwar einige meiner Top-Prioritäten, jedoch fehlen entscheidende Elemente wie C04 (508) und B85 (461), stattdessen sind Gegenstände mit sehr niedriger Priorität (A61:121, C49:262) enthalten. Ohne C04 ist mein Gesamtnutzen erheblich niedriger als möglich.'}\\ \tt VORSCHLAG: {'A48', 'A04', 'B21', 'B53', 'A41', 'C04', 'B85', 'C64'} 
	  } 
	   } 
	   } 
	 & \\ 
 

    \theutterance \stepcounter{utterance}  

    & & & \multicolumn{4}{p{0.6\linewidth}}{\cellcolor[rgb]{1,0.85,0.72}{%
	\makecell[{{p{\linewidth}}}]{% 
	  \tt {\tiny [GM$\langle$B]}  
	 STRATEGISCHE ÜBERLEGUNGEN: {'Derzeit wird der Vorschlag des Gegners (A48, A04, B21, B53, A41, C04, B85, C64) wiederholt. Meine optimale Strategie war zuvor das Hinzufügen von A61, C33 und C49, jedoch konnte aufgrund des Aufwandlimits kein Kompromiss gefunden werden. Da der Gegner nicht bereit ist, auf meine Prioritäten einzugehen, ist Akzeptanz des gegenwärtigen Vorschlags die beste Option, um den Verhandlungsprozess zu beenden und einen Nutzen von 3019 zu erzielen.'}\\ \tt ZUSTIMMUNG: {'A48', 'A04', 'B21', 'B53', 'A41', 'C04', 'B85', 'C64'} 
	  } 
	   } 
	   } 
	  \\ 
 

    \theutterance \stepcounter{utterance}  

    & & \multicolumn{4}{p{0.6\linewidth}}{\cellcolor[rgb]{0.9,0.9,0.9}{%
	\makecell[{{p{\linewidth}}}]{% 
	  \tt {\tiny [GM$\rangle$B]}  
	 Deine Antwort hat kein argument enthalten. Du muss das strukturierte Format ARGUMENT: {...} irgendwo in theiner Antwort verwenden. Versuche es erneut 
	  } 
	   } 
	   } 
	 & \\ 
 

    \theutterance \stepcounter{utterance}  

    & & & \multicolumn{4}{p{0.6\linewidth}}{\cellcolor[rgb]{1,0.85,0.72}{%
	\makecell[{{p{\linewidth}}}]{% 
	  \tt {\tiny [GM$\langle$B]}  
	 STRATEGISCHE ÜBERLEGUNGEN: {'Die aktuelle Verhandlung hat gezeigt, dass der Gegenspieler nicht bereit ist, auf meine höchsten Prioritäten einzugehen. Um den Prozess abzuschließen, akzeptiere ich den Vorschlag, da er zumindest die wichtigsten Elemente des Gegners enthält und der Gesamtaufwand im Limit liegt, was eine minimale akzeptable Lösung darstellt.'}\\ \tt ARGUMENT: {'Dieser Vorschlag berücksichtigt Ihre vorgegebene Konfiguration ohne weitere Restriktionen und ist somit der einzige verbleibende Kompromiss. Der Gesamtaufwand von 3921 erfüllt die Vorgaben und ermöglicht eine klare Lösung.'}\\ \tt ZUSTIMMUNG: {'A48', 'A04', 'B21', 'B53', 'A41', 'C04', 'B85', 'C64'} 
	  } 
	   } 
	   } 
	  \\ 
 

    \theutterance \stepcounter{utterance}  

    & & & \multicolumn{2}{p{0.3\linewidth}}{\cellcolor[rgb]{0.95,0.95,0.95}{%
	\makecell[{{p{\linewidth}}}]{% 
	  \tt {\tiny [GM$|$GM]}  
	 game successful 
	  } 
	   } 
	   } 
	 & & \\ 
 

    \theutterance \stepcounter{utterance}  

    & & & \multicolumn{2}{p{0.3\linewidth}}{\cellcolor[rgb]{0.95,0.95,0.95}{%
	\makecell[{{p{\linewidth}}}]{% 
	  \tt {\tiny [GM$|$GM]}  
	 end game 
	  } 
	   } 
	   } 
	 & & \\ 
 

\end{supertabular}
}

\end{document}
